% !TeX root = ../../Book5.tex
% 纲要标题：### 21.4 代数群结构初步
% 纲要位置：工作记录/计划纲要.md，第 4409 行。

\section{代数群结构初步}
\label{b5:sec:2104}


% 纲要范围：
%
% * 幺幂群和约化群
% * 根数据的意义
% * 正特征行为和群概形方向
%
% ---
%

从矩阵例子进入一般代数群时，需要区分几个不同的结构条件。幺幂性控制元素的 Jordan 形，约化性排除连通正规幺幂部分；根数据又在根系以外保留了整数格。正特征时，坐标环的幂零信息和无穷小子群还会使点集直觉失效。

\subsection{幺幂群和约化群}
\label{b5:subsec:2104:01}

在代数闭域 \(k\) 上，矩阵 \(u\) 若满足 \(u-I\) 幂零，就称为幺幂矩阵。一个线性代数群在一个忠实表示中全部元素都幺幂时，称为幺幂群；这个概念的表示无关性是一般线性代数群理论的一项结构结论。本节的实例均在明确的矩阵模型中使用它。
\ProseParagraphBreak
上三角对角线全为一的群 \(U_n\) 是基本例子。令 \(U_n^{(r)}\) 为仅在距对角线至少 \(r\) 条的上对角线上允许非零条目的子群。若记相应严格上三角矩阵空间为 \(J^r\)，则
\[
(1+X)(1+Y)\equiv1+X+Y\pmod{J^{r+s}}
\quad(X\in J^r,\ Y\in J^s),
\]
从而 \(\Bigl[U_n^{(r)},U_n^{(s)}\Bigr]\subseteq U_n^{(r+s)}\)。因为 \(J^n=0\)，这个群幂零，特别可解。对 \(U_3\)，三个参数的乘法为
\[
(a,b,c)(a',b',c')
=(a+a',b+b',c+c'+ab'),
\]
可见中心方向记录了非交换性，不能把全部上三角参数都简单看成独立加法。
\begin{definition}\label{b5:def:reductive-group}
设 \(G\) 为代数闭域上的光滑连通仿射代数群。若 \(G\) 没有非平凡的连通正规幺幂闭子群，就称为\term[yuehua-qun]{约化群}[reductive group]。若此外没有正维数的连通正规可解闭子群，就称为半单代数群。这里的“约化群”不同于“坐标环约化”：后者只要求没有非零幂零元。
\end{definition}
例如环面没有非平凡幺幂元素，所以约化；它自身交换且有正维数时不半单。\(U_n\) 在 \(n\ge2\) 时光滑连通且非平凡，却是幺幂的，因而不约化。一般约化群的结构定理还要研究环面、根子群和正规子群；本章只在下一小节解释其中不可省略的整数格资料，不在这里给出所有群的分类。

\subsection{根数据的意义}
\label{b5:subsec:2104:02}

\begin{definition}\label{b5:def:root-datum}
设 \(X,Y\) 为两个有限秩自由交换群，给定完美整数配对
\(\langle\ ,\ \rangle:X\times Y\to\mathbb Z\)，以及有限子集
\(\Phi\subseteq X,\Phi^\vee\subseteq Y\) 及双射
\(\alpha\mapsto\alpha^\vee\)。要求对所有 \(\alpha\in\Phi\)，都有
\(\langle\alpha,\alpha^\vee\rangle=2\)，且反射
\[
s_\alpha(x)=x-\langle x,\alpha^\vee\rangle\alpha,\qquad
s_\alpha^\vee(y)=y-\langle\alpha,y\rangle\alpha^\vee
\]
分别保持 \(\Phi\) 和 \(\Phi^\vee\)，并对所有 \(\alpha,\beta\in\Phi\) 满足 \((s_\alpha\beta)^\vee=s_\alpha^\vee(\beta^\vee)\)。还要求约化条件 \(\Phi\cap\mathbb Q\alpha=\{\alpha,-\alpha\}\) 对每个 \(\alpha\in\Phi\) 成立，其中把 \(X\) 嵌入 \(X\otimes_{\mathbb Z}\mathbb Q\)。满足这些条件的资料称为\term[gen-shuju]{根数据}[root datum]；允许 \(\Phi\) 为空，也不要求它生成整个 \(X\otimes\mathbb Q\)，以便包含环面的中心方向。
\end{definition}
在连通约化群的结构理论中，\(X\) 是极大环面的字符格，\(Y\) 是一参数子群格。根系给出实向量空间里的角度与相对长度，根数据还记录哪些整数权给出群上的字符。因此同一个复半单李代数可以对应不同的代数群。
\ProseParagraphBreak
对 \(\operatorname{SL}_2\)，取 \(T=\Bigl\{\operatorname{diag}(t,t^{-1})\Bigr\}\)。若 \(\lambda(t)=t\)，则
\(X=\mathbb Z\lambda\)，根为 \(\pm2\lambda\)；一参数子群
\(\nu(t)=\operatorname{diag}(t,t^{-1})\)
满足 \(\langle\lambda,\nu\rangle=1\)，正余根就是 \(\nu\)。
对 \(\operatorname{PGL}_2\) 的对角环面，用对角条目比值 \(r\) 作坐标，字符格由 \(\alpha(r)=r\) 生成，根为 \(\pm\alpha\)，余根是相应余字符生成元的二倍。两者的抽象根系都为 \(A_1\)，但整数格中的嵌入不同。
\ProseParagraphBreak
这也解释了表示下降的一个必要条件。中心元 \(-I\in\operatorname{SL}_2(\mathbb C)\) 在
\(L_m=\operatorname{Sym}^m(\mathbb C^2)\)
上作用为 \((-1)^m\)，所以这个具体代数群表示能通过商
\(\operatorname{PSL}_2(\mathbb C)\) 当且仅当 \(m\) 为偶数。Lie 代数的最高权分类列出所有 \(m\ge0\)，却不会自动替每个中心商筛去不能下降的表示。

\subsection{正特征下的代数群与群概形}
\label{b5:subsec:2104:03}

在特征 \(p\) 中，对加法群的 Frobenius 态射
\[
F:\mathbb G_a\longrightarrow\mathbb G_a,\qquad t\longmapsto t^p
\]
对应坐标同态 \(k[t]\to k[t]\)、\(t\mapsto t^p\)。它的核为
\(\alpha_p=\operatorname{Spec}k[t]/(t^p)\)。
若 \(k\) 代数闭，点集上 \(F\) 为双射，但坐标同态不满，所以它不是代数群同构。其微分为零；核概形的切空间却一维，尽管它的 \(k\) 点只有一个。
\ProseParagraphBreak
同样，\(\mu_p=\operatorname{Spec}k[t,t^{-1}]/(t^p-1)\) 的切空间一维；在 \(t=1+\epsilon a\) 代入关系时，
\((1+\epsilon a)^p=1\)，对任意 \(a\) 都成立。于是 \(\alpha_p,\mu_p\) 与平凡点集之间的区别，必须由坐标环或所有代数上的点保留。
\ProseParagraphBreak
这些例子不是给特征零理论添上的小修饰。它们说明光滑性、约化性、可分性以及表示是否能由微分恢复，分别控制不同信息。后续的不变量理论同样会明确区分特征零线性约化性与正特征下更弱的几何约化性，不能用有限群的平均公式替代一般群概形的论证。
